% Part: first-order-logic
% Chapter: sequent-calculus

\documentclass[../../../include/open-logic-chapter]{subfiles}

\begin{document}

\iftag{FOL}
      {\olchapter{fol}{seq}{సీక్వెంట్ కలనం}}
      {\olchapter{pl}{seq}{సీక్వెంట్ కలనం}}

\begin{editorial}
  ఈ అధ్యాయం సాంప్రదాయిక మొదటిస్థాయి విధేయ తర్కానికి గెంట్సెన్ ప్రతిపాదించిన
  ప్రామాణిక సీక్వెంట్ కలనం \Log{LK}ను పరిచయం చేస్తుంది. దీనికి మరిన్ని
  ఉదాహరణలు, అభ్యాసాలు అవసరం. నిరూపణ వ్యవస్థగా సీక్వెంట్ కలనానికి
  సంబంధించిన విషయాన్ని చేర్చడానికి లేదా మినహాయించడానికి ``prfLK''
  ట్యాగును ఉపయోగించండి.
\end{editorial}

\olimport{rules-and-proofs}

\olimport{propositional-rules}

\iftag{FOL}{%
  \olimport{quantifier-rules}
}{}

\olimport{structural-rules}

\olimport{derivations}

\olimport{proving-things}

\iftag{FOL}{%
  \olimport{proving-things-quant}
}{}

\olimport{proof-theoretic-notions}

\olimport{provability-consistency}

\olimport{provability-propositional}

\iftag{FOL}{%
  \olimport{provability-quantifiers}
}{}

\olimport{soundness}

\iftag{FOL}{%
  \olimport{identity}
  \olimport{soundness-identity}
}{}

\OLEndChapterHook

\end{document}
